\section{Fourier transforms on prime tensor powers}
\label{sec:prime-kernel}

On \(k\) digits of radix \(\ell\), the Fourier matrix is the tensor
product of \(k\) copies of a single \(\ell\times\ell\) matrix. Applying
the factors in succession takes \(O_\ell(k\ell^k)\) operations. We
reduce the power of \(k\) by applying the prediction and repair of
Theorem~\ref{thm:repair} to a batch of independent transforms.

For a prime \(\ell\), let \(\zeta=\exp(2\pi i/\ell)\), put
\(R_\ell=\mathbb Z[1/\ell,\zeta]\), and write
\(D_\ell=(\zeta^{xy})_{x,y\in\mathbb F_\ell}\). The rows and columns
of \(D_\ell^{\otimes k}\) are indexed by \(\mathbb F_\ell^k\).

\begin{theorem}[Fourier transforms in a fixed prime radix]\label{thm:kernel}
For each prime \(\ell\) and every \(\alpha>1/2\), there are effectively
constructible \(R_\ell\)-linear circuits for \(D_\ell^{\otimes k}\)
of size \(O_{\ell,\alpha}(\ell^k(1+k)^\alpha)\).
The same bound counts exact ring operations, including scalar
multiplications, and indexing on \(O(1+k)\)-bit address words.
Integer arithmetic, comparisons, lookup and digit extraction on a
word have unit cost. For fixed \(\ell,\alpha\) the control depends only
on \(k\), and the coefficient identities specialize under every
homomorphism from \(R_\ell\).
\end{theorem}

Group the digits into \(A\) axes of \(f\) digits each. A \emph{call}
applies \(D_\ell^{\otimes f}\) along one axis of an array of
\(\ell^{Af}\) entries, at every choice of the other digits. The
usual tensor product algorithm uses \(A\) calls per array. Fix
\(1/2<\theta\le\alpha\). Our goal at one recursion level is to
transform \(W\) independent arrays in \(\sigma\) calls with
\(\sigma/W<A^\theta\). The integers \(A,W,\sigma\) will depend on
\(\ell,\theta\) and work for every \(f\). The strict inequality leaves
room to pad incomplete batches during recursion.

The large transform will be a product of \(m\) commuting queries
up to a permutation and diagonal scalings (together called a
\emph{monomial map}). The number of axes \(A\) will be of order \(m\).
Theorem~\ref{thm:repair} replaces the product by
\(O_\ell(\sqrt{m\log m})\) pencil stages and full Fourier
conversions on a small fraction of the rows. A pencil still contains all \(m\) queries. The construction
below makes their Fourier calls shareable across queries and across the
independent arrays in the batch.

\subsection{The product of queries}

Choose a large \(m\equiv1\pmod\ell\). In each half of an address use
\(m\) distinguished axes, one per query, and \(D-1\) auxiliary axes
shared by all the queries. The routing will determine \(D\). Put
\(d=m+D-1\) and \(A=2d\). Read the address by columns. Its column
\(c\) contains the \(c\)-th digit of each axis, hence a pair of
\(d\)-vectors. Write \(X_c,Z_c\) for the frequency coordinates
in column \(c\). The \(j\)-th query, \(1\le j\le m\), has multiplier
\[
 Q_j(X,Z)=\zeta^{\sum_{c=1}^{f}
 (X_{c,j}Z_{c,j}+\sum_{i=m+1}^dX_{c,i}Z_{c,i})}.
\]
More precisely, \(Q_j(X,Z)\) is the diagonal entry of
\(\mathcal S Q_j\mathcal S^{-1}\), where
\(\mathcal S=D_\ell^{\otimes Af}\).
Each distinguished coordinate occurs once in the product
\(F=\prod_j Q_j\), and each auxiliary coordinate occurs \(m\equiv1
\pmod\ell\) times. Thus its exponent is \(\sum_c X_c\cdot Z_c\).
Fourier inversion for one digit pair gives the ordinary kernel
from \((x,y)\) to \((s,t')\),
\[
 \ell^{-2}\sum_{X,Z}\zeta^{(x-s)X+(y-t')Z+XZ}
 =\ell^{-1}\zeta^{(x-s)(t'-y)}
 =\ell^{-1}\zeta^{-xy-st'}\zeta^{xt'+ys}.
\]
Consequently \(F\) gives \(\ell^{-df}D_\ell^{\otimes Af}\) up to
diagonal phases and exchange of the two halves. The scalar
\(\ell^{df}\) recovers the unnormalized transform.

Apply Theorem~\ref{thm:repair} on \(R\) independent rows. A pencil
has terms \(B_{ja}\otimes Q_j^a\), \(1\le a<\ell\), and a pointwise
term. The matrix \(B_{ja}\) acts on the \(R\) rows. Across all stages
there are \(O_\ell(m^{-3}R)\) full-row spectral conversions. The
coefficients, row count and spectral tables are finite once \(m\) is fixed.

\subsection{Sharing the queries of a pencil}

Run the \(R\)-row program in parallel at each label \(v\) of a finite
set \(\mathcal V\). The desired transfer on labels is the identity,
so a source must contribute its query to its own destination only.
To realize the identity, we route data through weighted rectangles of
labels. A rectangle \(I\times J\) sums the arrays from sources in \(I\)
and sends the result to destinations in \(J\).

Fix a query \(j\) and a power \(a\). At each label \(v\), choose
input and output permutations \(S_{v,j}\) and \(T_{v,j}\) into a
common coordinate system. We shall construct a kernel \(\mathcal K_a\)
in these coordinates with
\(T_{v,j}^{-1}\mathcal K_a S_{v,j}=Q_j^a\).
Routing with the label matrix \(M\) gives the operator block
\begin{equation}\label{eq:kernel-routing-block}
  M_{wv}\,T_{w,j}^{-1}\mathcal K_a S_{v,j}
  \qquad\text{from source \(v\) to destination \(w\)} .
\end{equation}
Thus summing rectangle matrices to the identity cancels all
off-diagonal transfers, even though labels use different permutations.
The rectangles have two different costs. Some require a common
Fourier computation for each query. On the others, the distinguished
directions of every source and destination are orthogonal. This lets
us reuse the same endpoint Fourier computations for all \(j\).

Following the intersection representations of
\cite{Alon1998,Tape}, we use one matrix over \(\mathbb F_\ell\) to
choose directions and a second matrix over \(\mathbb Z[1/\ell]\) to
choose rectangle weights. The rank of the first determines the number
of auxiliary axes. The second bounds the number of rectangles charged
separately for every query.

Fix an integer \(L\ge\ell\) and a power of two \(q\). The label set
\(\mathcal V=\{0,\ldots,q-1\}^L\) has cardinality \(q^L\). Write
\(\Delta(v,w)\) for Hamming distance and define the symmetric matrix
\[
 \Gamma(v,w)=1-\Delta(v,w)^{\ell-1}\quad\hbox{in }\mathbb F_\ell .
\]
Its rank \(D\) is \(O_{\ell,L}(q^{\ell-1})\). Indeed, expand in the
indicators of matches in the \(L\) positions. A product of match
indicators on a set of \(j\le\ell-1\) positions has rank at most
\(q^j\), by summing the indicator outer products for the common
subwords. Give \(\mathbb F_\ell^{\mathcal V}\) the form with Gram
matrix \(\Gamma\) and take its quotient by the radical. This gives a
\(D\)-dimensional space \(\mathcal D\) with a nondegenerate symmetric
pairing and vectors \(d_v\) such that
\(\langle d_v,d_w\rangle=\Gamma(v,w)\). In particular
\(\langle d_v,d_v\rangle=1\).

We now decompose the identity on \(\mathcal V\) into rectangles over
\(\mathbb Z[1/\ell]\). All but a small number will satisfy
\(\langle d_v,d_w\rangle=0\) throughout the rectangle.
Put \(t=\lfloor L/\ell\rfloor\) and
\[
 p(k)=\prod_{h=1}^{t}\left(1-\frac{k}{\ell h}\right).
\]
For an integer \(k\), \(p(k)\) belongs to \(\mathbb Z[1/\ell]\).
For every prime \(b\ne\ell\) and \(e\ge1\), at least
\(\lfloor t/b^e\rfloor\) of the factors \(\ell h-k\) are divisible
by \(b^e\), since \(h\) runs through consecutive residues and \(\ell\)
is invertible modulo \(b^e\). This accounts for the \(b\)-part of
\(t!\). Expand \(p(k)\) in the basis \(\binom{L-k}{j}\), \(0\le j\le t\).
Evaluation at \(k=L,L-1,\ldots,L-t\) gives a unit triangular system,
so its coefficients belong to the same ring. At \(k=\Delta(v,w)\)
the \(j\)-th binomial counts \(j\)-subsets of equal positions. Fixing
both the subset and its common letters gives in all
\(O_L(q^t)\) rectangles. We charge these rectangles separately for
each query.

The difference between the identity and \(p(\Delta)\) is supported
at distances nonzero modulo \(\ell\), since \(p(0)=1\) and
\(p(\ell h)=0\) for \(1\le h\le t\). At each such distance fix the
set of unequal positions. Fix the common letters at the other
positions. At each unequal position partition the ordered pairs of
letters by their first differing bit, their common preceding prefix,
and the order of the two differing bits. Each part is a rectangle.
These products of rectangles, weighted by \(-p(\Delta)\), complete
the identity. On each of them \(\Gamma|_{I\times J}=0\). We call them
free rectangles because their Fourier operations can be shared over
all the queries. For each label, the number of incidences on either
side, including the charged rectangles, is
\(O_L((1+\log q)^L)\). Indeed, at a fixed source or destination letter there is one equality
choice and at most one choice for each possible first differing bit.

\subsection{The endpoint Fourier transforms}

For fixed \(j\), use in each half the distinguished coordinate \(j\)
and the \(D-1\) auxiliary coordinates \(m+1,\ldots,d\). In the
ordinary array coordinates of label \(v\), identify both resulting
\(D\)-vectors with \(\mathcal D\). Use \(d_v\) in the distinguished
direction, and a basis and its paired dual in \(d_v^\perp\) for the
two lists of auxiliary coordinates. At each label, both bases are fixed independently of \(j\).
They exist because \(\langle d_v,d_v\rangle=1\). Use these
identifications on each of the \(f\) columns on input and output, and
keep every spectator axis in its original order. The auxiliary digits
belong to the address of the input array and carry data throughout.
This defines \(S_{v,j}\) and \(T_{v,j}\) in
Equation~\eqref{eq:kernel-routing-block}.

Since the two chosen bases are dual,
the original dot product between the first and second halves becomes
the pairing on \(\mathcal D\). In common coordinates \(Q_j^a\) has
one-column kernel
\[
 K_a(s,t';x,y)=\ell^{-D}\,
              \zeta^{\langle x-s,t'-y\rangle/a},
 \qquad x,y,s,t'\in\mathcal D,
\]
by the same summation with \(aXZ\) in the exponent. The full
\(\mathcal K_a\) is its product over the columns, tensored with the
identity on the unused axes. The local spectator digits depend on \(j\).

For a charged rectangle the common-coordinate sum has \(R\) rows.
Apply \(\mathcal K_a\) to each row, distribute it to the destinations
with the rectangle weight, and apply \(B_{ja}\) across row indices.
This costs at most \(2DR\) calls at the specified \(j,a\). The
cross terms use \(D\) one-axis calls on each half, and the remaining
factors are monomial maps.

For a free rectangle \(I\times J\), factor the common kernel into
two partial transforms as in Figure~\ref{fig:route}. Set
\(H=\operatorname{span}\{d_w:w\in J\}\) in common coordinates,
and choose a vector space complement \(C\) to \(H\) in
\(\mathcal D\). For the ordinary input
coordinates write \(x=(x_H,x_C)\), \(y=(y_H,y_C)\) in a basis
adapted to \(H,C\). Write \(s=(s_1,s_2)\), \(t'=(t'_1,t'_2)\) in
the paired dual basis. After extracting the initial and final
diagonals \(\zeta^{-\langle x,y\rangle/a}\) and
\(\zeta^{-\langle s,t'\rangle/a}\), the cross term is
\[
 \zeta^{(x_H\cdot t'_1+y_H\cdot s_1)/a}\,
 \zeta^{(x_C\cdot t'_2+y_C\cdot s_2)/a}.
\]
The first partial transform sums along \(H\) in each input at
fixed \(x_C,y_C\) (and fixed values on spectator axes). Its intermediate coordinates are
\((x_C,y_C,s_1,t'_1)\). The second partial transform produces
\((s_2,t'_2)\) at fixed \(s_1,t'_1\). In the transposed kernel this
second transform sums along \(H^\perp\) in the output, since the
vectors dual to the \(C\) part span \(H^\perp\). This argument uses a
basis and its dual for the full nondegenerate pairing, and also
applies when the restricted form on \(H\) is degenerate.

To share the first transform over \(j\), express it in the coordinates
of a source \(v\in I\). Since \(H\subseteq d_v^\perp\), a coset of
\(H\) can be represented using the distinguished coordinate and a
complement to \(H\) inside \(d_v^\perp\). Summation along \(H\)
then uses only auxiliary axes. Switching to the common representative
from \(C\) permutes the cosets and scales the Fourier coefficients.
For example, replacing a representative of \(x+H\) by that
representative plus \(h\in H\) scales the coefficient at \(t'_1\)
by \(\zeta^{-h\cdot t'_1/a}\). This remains true with the initial
diagonal phase included in the input. That phase factors into an
auxiliary phase and a distinguished phase, since the two directions
are orthogonal. The distinguished phase is constant on each coset and
can be applied after the transform. What remains at the source is
one auxiliary transform independent of \(j\).

At a destination \(w\in J\) the analogous statement follows from
\(H^\perp\subseteq d_w^\perp\). In the transposed kernel the final
transform runs along \(H^\perp\), hence also on auxiliary axes only.
Between these two endpoint computations, each \(j\) uses just
monomial changes of coordinates and pointwise mixing. The row matrix
\(B_{ja}\) commutes with the final transform, so the destination can
sum over \(j\) before taking that transform. All partial transforms
and monomial changes act column by column. Mixing by \(B_{ja}\) acts
across rows at each address. These operations therefore implement
Equation~\eqref{eq:kernel-routing-block} on the full \(Af\)-digit arrays.

\begin{figure}[tbp]
\centering
\begin{tikzpicture}[
  x=1cm,y=1cm,font=\small,align=center,
  endpoint/.style={draw=#1,fill=#1!5,line width=.55pt,rounded corners=2pt,
        text width=2.55cm,minimum height=2.25cm,inner xsep=3pt},
  query/.style={draw=ink!70,fill=ink!3,rounded corners=2pt,
        text width=4.3cm,inner sep=4pt},
  edge/.style={draw=ink,-{Stealth[length=1.6mm]},line width=.48pt}]
\node[endpoint=sea] (source) at (-4.3,0)
  {\textbf{Source \(v\in I\)}\\[3pt] Partial Fourier transform along \(H\)};
\node[endpoint=rust] (dest) at (4.3,0)
  {\textbf{Destination \(w\in J\)}\\[3pt] Complete the transform along the dual axes};
\node[query] (first) at (0,.85)
  {Reindex, scale and mix rows by \(B_{1a}\)};
\node[font=\small] at (0,0) {\(\vdots\)};
\node[query] (last) at (0,-.85)
  {Reindex, scale and mix rows by \(B_{ma}\)};
\foreach \op/\off in {first/5mm,last/-5mm}{
  \draw[edge] ([yshift=\off]source.east)--(\op.west);
  \draw[edge] (\op.east)--([yshift=\off]dest.west);
}
\node[font=\footnotesize,text=ink] at (0,1.7)
  {Sum source arrays over \(I\) in common coordinates, then copy to \(J\).};
\node[font=\footnotesize,text=ink] at (0,-1.7)
  {At each destination, sum the \(m\) contributions before its Fourier transform.};
\end{tikzpicture}
\caption{Fourier calls for one free rectangle \(I\times J\) and one power
\(a\). A source takes its partial transform once for all \(m\) queries.
Between the endpoints the common coordinates are \((x_C,y_C,s_1,t'_1)\)
in each column. The transposed destination transform acts along
\(H^\perp\). Both endpoint transforms use only auxiliary axes.}
\label{fig:route}
\end{figure}


\subsection{The finite saving and recursion}

Each endpoint uses at most \(2D\) calls per row incident to the
rectangle and per \(a\). For charged rectangles there are
\(O_L(q^t)\) common transforms per \(j,a\) across all \(q^L\) labels.
The number of calls in a pencil divided by the total \(Rq^L\) row
count is therefore at most
\begin{equation}
 \label{eq:kernel-pencil-cost}
 C_{\ell,L}D\bigl((1+\log q)^L + m q^{-(L-t)}\bigr).
\end{equation}
One full-row spectral conversion uses \(2d\) calls, also for the
inverse by frequency negation and scalar normalization. The repair
conversions therefore contribute \(O_\ell(dm^{-3})\) calls per row.
Choose \(q\) to be a power of two within a factor two of
\(m^{1/(L-t)}\). Then \(m q^{-(L-t)}=O_L(1)\).
The dimension estimate gives
\(D=O_{\ell,L}(m^{(\ell-1)/(L-t)})\). Given
\(0<\eta<\min(1,\theta-1/2)\), choose \(L\) large enough to make
\((\ell-1)/(L-t)<\eta\). The factor \((1+\log q)^L\) is then
absorbed by the remaining power of \(m\), so
Equation~\eqref{eq:kernel-pencil-cost} is \(O_{\ell,L}(m^\eta)\)
and \(D=o(m)\). Over all stages, the call count per row is at most
\[
 O_{\ell,L}\bigl(\sqrt{m\log m}\,m^\eta+dm^{-3}\bigr)
       =o(A^\theta).
\]
Here \(A=2(m+D-1)\). We may now fix \(m\equiv 1\pmod\ell\) large
enough to obtain the required strict saving at one recursion level.

The coefficients of the finite repair may lie in a fixed finite
free extension of \(R_\ell\). Expand each row in a basis of the
extension. A pointwise multiplication becomes a fixed matrix over
\(R_\ell\). Every axis call acts separately on basis components.
The target transform has base coefficients and the program accepts
arbitrary extension-valued input. Hence each basis component carries
an independent array over \(R_\ell\) and receives its own transform.
Expanding the scalars multiplies both the call count and the number
of input arrays by the extension rank. If \(W\) is the expanded number of rows and
\(\sigma\) the expanded number of calls, we have proved
\begin{equation}
 \label{eq:kernel-template-saving}
                          \sigma/W < A^\theta .
\end{equation}
For every \(f\ge1\), a batch contains \(W\) input and \(W\) output
rows of \(\ell^{Af}\) entries each. Each call acts on a single work
row of \(\ell^{Af}\) entries, applying \(D_\ell^{\otimes f}\) along
one axis for all \(\ell^{(A-1)f}\) fibers. An arbitrary common set of
spectator coordinates multiplies all these volumes equally. Outside
calls there are a fixed number of work rows and of linear passes.

Fix \(r\) so large that
\(b=(\sigma/W)(1+W/\ell^r)<A^\theta\). For a single \(k\)-digit array
take \(f=\lfloor(k-r)/A\rfloor\) when positive. The \(k-Af\) inactive
digits provide independent rows for the finite template. Batch these
rows with padding at the last batch. Handle the bounded number
\(k-Af<r+A\) of remaining digit transforms directly. The total work
divided by \(\ell^k\) satisfies
\[
                         T(k)\le C+bT(f).
\]
Indeed, each call runs the single-array algorithm independently on
its fibers and covers one full work row. The padding factor satisfies
\(\lceil\ell^{k-Af}/W\rceil W/\ell^{k-Af}\le 1+W/\ell^r\).
For bounded \(k\) use direct transforms. The ratio
\(b(1+f)^\theta/(1+k)^\theta\) has limit \(b/A^\theta<1\), so
induction gives \(T(k)=O((1+k)^\theta)\).

For uniformity, store the \(f\) columns in positional layout.
All index changes in the common-coordinate and representative
conversions apply a fixed linear change to each column. Interchanging
axes and changing strides in packing fibers also have known positional
weights. Enumerate column prefixes, extending each partial address
and its image by the next digit tuple times these weights. There are
a constant number of extensions per prefix and \(O(\ell^{Af})\)
prefixes. Along the same enumeration accumulate the phase exponents
and the \(m\) query exponents modulo \(\ell\). The query exponents
specify the repair table. All fixed ring constants have fixed expressions in
\(\zeta,1/\ell\) with integer coefficients. Normalization powers
depending on \(f\) take \(O(1+\log f)\) ring operations. These
procedures give linear work in local volume plus polynomial setup
in the local exponent, absorbed in that volume. Packing and
unpacking each batch requires a constant number of such traversals.
Address words of \(O(k+1)\) bits suffice, since the total number of
gates and temporary addresses is polynomial in \(\ell^k\). This
proves Theorem~\ref{thm:kernel}.


\subsection{Bit layout and stability for the Walsh--Hadamard transform}
\label{sec:binary}

At radix two, each recursive call transforms selected bits of an
address and leaves the other bits available for independent inputs.
The tape multiplication algorithm in \cite{oaT} uses this layout.
The same template also computes the normalized Walsh--Hadamard
transform. Its normalization can be absorbed into the recursive
calls, giving a fixed coefficient set and a bound on the accumulation
of numerical error.

\begin{proposition}[Binary implementations]\label{prop:binary-template}
Fix \(\theta>1/2\) and the arity \(A\), batch size \(W\) and call count
\(\sigma\) of Equation~\eqref{eq:kernel-template-saving} at \(\ell=2\).
\begin{enumerate}
\item An unnormalized batch on \(Af\) bits uses, outside its recursive
calls, pointwise linear maps on a fixed number of row entries, row
rewirings, and invertible \(\mathbb F_2\)-linear reindexings of the
\(A\) bits within each column. Its coefficients are drawn from a fixed
finite list of dyadic rationals and from signed powers of two with
exponent of magnitude \(O(f+1)\). Each call transforms a single
\(f\)-bit axis independently at every setting of the other bits.

The \(A\) axes may occupy contiguous slots consisting of \(f\) blocks
of \(K\) bits each. For any fixed \(0\le\rho<K\), the template can use
just bit \(\rho\) of each block. All other bits retain their addresses.
\item For even \(k\), the normalized transform \(2^{-k/2}D_2^{\otimes k}\)
has a uniform dyadic linear circuit of size
\(O_\theta(2^k(1+k)^\theta)\) with coefficients in a fixed finite
list outside the recursive calls. There are constants \(B,C,c>0\)
depending only on the template with the following property. If each
gate value is multiplied by \(1+\delta_g\), where
\(\lvert\delta_g\rvert\le u\le c(1+k)^{-B}\), the perturbed and exact
maps differ in Euclidean operator norm by at most \(Cu(1+k)^B\).
\end{enumerate}
\end{proposition}

\begin{proof}
The coefficient ring \(R_2=\mathbb Z[1/2]\) makes the fixed row
operations dyadic after descent. Table lookups use sums modulo two
of fixed column functions, and the phases are signs. The remaining
coefficients are powers of two from the axis normalizations and
from recovering the unnormalized target from the product \(F\).

By binary elimination each column reindexing is a fixed sequence of
additions from one axis to another, with an exchange implemented by
three additions. If \(x\) and \(y\) are the source and destination
slots for one such addition, the update is
\begin{equation}\label{eq:reindex}
 y_{\rho+iK}\longleftarrow y_{\rho+iK}+x_{\rho+iK}
       \quad\text{in }\mathbb F_2,\qquad 0\le i<f .
\end{equation}
Only the selected bits change. Fixing \(s\) inactive bits divides
one array of \(Z\) entries into \(2^s\) independent streams of
\(Z/2^s\) entries each. Group these streams into batches of \(W\),
padding the last batch. The recursive calls process in total
\[
 \sigma\lceil2^s/W\rceil\, Z/2^s
       \ \le\ (\sigma/W)(1+W/2^s) Z
\]
entries. This bound continues to hold for every slot width and offset.

For the normalized transform, take \(f\) to be the largest even
integer at most \((k-r)/A\) whenever this yields a positive child
width. The remaining \(k-Af<r+2A\) bits have even width and can be
handled by normalized four-point butterflies, as can the bounded
base cases. The padding bound and the strict saving in the recurrence
still apply.

Use a normalized transform in each axis call. For a common kernel
the \(2D\) calls supply the entire factor
\((2^{-f/2})^{2D}=2^{-Df}\). For a free rectangle the two partial
transforms together use those same \(2D\) axes. In a full spectral
conversion the common normalization commutes through its row table
and cancels with that of the inverse. Finally, \(F\) is already the
normalized target up to phases and a permutation. Thus all scalars
outside the recursive calls and the bounded remainder belong to a
fixed finite dyadic list.

Give the concatenated work rows their Euclidean norm. Each pointwise
stage is a direct sum of fixed-size circuits with bounded coefficients,
so its exact operator norm is bounded and perturbing its gates changes
the stage by \(O(u)\) for \(u\le1\). Permutations, zero padding and
restriction have norm at most one. Each recursive stage acts on
independent fibers by unitary transforms. Its perturbation has norm
bounded by the largest error on a fiber, even if different gates have
different errors. The number of successive stages per recursion level
is fixed. Telescoping their product therefore bounds the worst error
\(e(k)\) by \(C_0(u+e(f))\) whenever \(e(f)\le1\), with a constant
that includes the bounded remainder and base cases. The recursion
has depth \(O(1+\log_A(1+k))\). Iteration gives
\(e(k)\le Cu(1+k)^B\), and choosing \(c\) sufficiently small ensures
the child bound \(e(f)\le1\) at every level.
\end{proof}

Over \(\mathbb Q(\sqrt2)\) one can normalize odd widths as well.
Remove one bit from an odd-width call, apply \(2^{-1/2}D_2\) to that
bit, and use the even-width circuit on the rest. One may equally
normalize every axis call in a one-level template with arbitrary \(f\),
since its \(2D\) factors still multiply to \(2^{-Df}\). This retains
the call count, the inactive coordinates and the columnwise
reindexings. The scalars outside the calls now form a fixed finite
algebraic list. They can be approximated to \(p\) binary places in
time polynomial in \(p\) using dyadic arithmetic and the positive
root of \(2x^2=1\).

